% SPDX-License-Identifier: CC-BY-4.0
\section{The upper bound}\label{sec:upper}

From now on the field is \(\Ft\) unless stated otherwise.

\subsection{Binary digit sums}

Let \(s_2(m)\) be the sum of the binary digits of \(m\ge0\).

\begin{lemma}\label{lem:digits}
For all \(a,b,j,J\ge0\):
\begin{enumerate}
 \item \(s_2(a+b)\le s_2(a)+s_2(b)\) and \(s_2(2^j)=1\). Hence a sum of \(r\)
 powers of two has digit sum at most \(r\).
 \item \(s_2(a)\le a\).
 \item If \(a\le b\), then \(2a+s_2(b)\le2b+s_2(a)\). That is,
 \(a\mapsto2a-s_2(a)\) is monotone.
 \item \(s_2(2^J-1)=J\).
\end{enumerate}
\textnormal{\small\leanref{claims/TruncatedProduct.lean}{[Lean: parts 1--3]}\,\leanref{claims/BinaryMultiplicityDegreeFormula.lean}{[Lean: part 4]}}
\end{lemma}

\begin{proof}
All four parts use the recursion \(s_2(m)=(m\bmod2)+s_2(\lfloor m/2\rfloor)\).
\begin{enumerate}
 \item The first inequality is proved by strong induction on \(a+b\), using
 \(\lfloor(a+b)/2\rfloor=\lfloor a/2\rfloor+\lfloor b/2\rfloor+c\), where
 \(c=\lfloor((a\bmod2)+(b\bmod2))/2\rfloor\) is the carry. The value
 \(s_2(2^j)=1\) follows from the recursion by induction on \(j\).
 \item This is standard; Mathlib proves it for every base.
 \item Induct on \(b-a\), using \(s_2(b+1)\le s_2(b)+1\), which is part~1 with
 second summand \(1\).
 \item \(2^{J+1}-1\) is odd with \(\lfloor(2^{J+1}-1)/2\rfloor=2^J-1\). \qedhere
\end{enumerate}
\end{proof}

\begin{lemma}[Legendre form]\label{lem:legendre}
Let \(n\ge0\), \(q\ge0\) and \(0\le r<2^n\), and put \(m=q2^n+r\). Then
\[
 \sum_{j=0}^{n-1}\Bigl\lfloor\frac m{2^j}\Bigr\rfloor+2q+s_2(r)=2m .
\]
Hence, for \(n\ge1\), \(0\le\ell<k\) and \(k-\ell-1=q2^n+r\) with \(0\le r<2^n\),
\(\Phi(n,k,\ell)=n+2k-2-2q-s_2(r)\), and \(\Phi(n,k,\ell)=n+2k-2-s_2(k-\ell-1)\) when
\(k-\ell-1<2^n\). Moreover \(m\) is the sum of \(h=2q+s_2(r)\) powers of two, each
at most \(2^{n-1}\) when \(n\ge1\): \(2q\) copies of \(2^{n-1}\) and the binary
digits of \(r\)
\textnormal{\small\leanref{claims/PerOrderLegendre.lean}{[Lean]}\,\leanref{claims/PerOrderArithmetic.lean}{[Lean: form of \(\Phi\)]}}.
\end{lemma}

\begin{proof}
For \(j<n\), \(\lfloor m/2^j\rfloor=q2^{n-j}+\lfloor r/2^j\rfloor\), and
\(\sum_{j<n}q2^{n-j}=2q2^n-2q\). Since \(r<2^n\), \(\sum_{j<n}\lfloor r/2^j\rfloor
=\sum_{j\ge0}\lfloor r/2^j\rfloor\), which is \(2r-s_2(r)\). This is Legendre's
classical identity \(\sum_{j\ge1}\lfloor r/2^j\rfloor=r-s_2(r)\), plus the term
\(j=0\); by induction, for \(r=2r'+b\) with
\(b\in\{0,1\}\) the sum is \(r+(2r'-s_2(r'))=2r-s_2(r)\). Adding gives
\(2m-2q-s_2(r)\). The formula for \(\Phi\) follows with \(m=k-\ell-1\), since
\(n+2\ell+2m=n+2k-2\). For the decomposition, \(2q\cdot2^{n-1}=q2^n\), and the
binary digits of \(r<2^n\) are powers \(2^j\) with \(j\le n-1\).
\end{proof}

\subsection{The truncated product}

Put \(y_i=x_i^2+x_i=x_i(1+x_i)\).

\begin{lemma}\label{lem:telescope}
In every commutative ring of characteristic two and for every \(t\) and \(K\ge0\),
\(\sum_{j<K}(t^2+t)^{2^j}=t^{2^K}+t\). In particular, over \(\Ft\),
\[
 (1+x)+\sum_{j<K}y^{2^j}=(1+x)^{2^K},\qquad y=x^2+x
\]
\lean{claims/TruncatedProduct.lean}.
\end{lemma}

\begin{proof}
By Frobenius, \((t^2+t)^{2^j}=t^{2^{j+1}}+t^{2^j}\), so the sum telescopes. For the
second form take \(t=1+x\), for which \(t^2+t=x^2+x\).
\end{proof}

\begin{definition}[Truncated product]\label{def:truncated}
For \(s\ge1\), let \(g_s\in\Ft[x_\sigma]\) be the sum, over all subsets
\(S\subseteq\sigma\) and all assignments \(i\mapsto p_i\) of powers of two to
\(i\in S\) with \(\sum_{i\in S}p_i\le s-1\), of
\[
 \prod_{i\in S}y_i^{p_i}\prod_{i\notin S}(1+x_i)
\]
\lean{claims/TruncatedProduct.lean}.
\end{definition}

\begin{lemma}[Catalan parity {\cite{AK73}, \cite[Theorem~2.1]{DS06}}]\label{lem:catalan}
For \(a\ge1\), \(v_2(C_{a-1})=s_2(a)-1\). In particular the Catalan number
\(C_{a-1}\) is odd if and only if \(a\) is a power of two
\lean{third-party-claims/CatalanParity.lean}.
\end{lemma}

\begin{proof}
By Kummer's theorem, \(v_2\binom{2n}{n}\) is the number of carries when adding \(n\)
to itself in base two, which is \(s_2(n)\). Since \((n+1)C_n=\binom{2n}{n}\),
\(v_2(C_n)=s_2(n)-v_2(n+1)=s_2(n+1)-1\), where the last step is
\(s_2(n+1)+v_2(n+1)=s_2(n)+1\): adding one clears the \(v_2(n+1)\) trailing ones of
\(n\) and sets one bit. Take \(n=a-1\).
\end{proof}

For \(s\ge2\), Menezes' Catalan truncation~\cite[equation~(12)]{Men26} is the
integral polynomial
\[
 \widehat g_s=(-1)^{s-1}\sum_{S,\,(a_i)}(-1)^{\sum_{i\in S}a_i}
 \prod_{i\in S}C_{a_i-1}\bigl(x_i(x_i-1)\bigr)^{a_i}\prod_{j\notin S}(x_j-1),
\]
summed over \(S\subseteq\sigma\) and integers \(a_i\ge1\) (\(i\in S\)) with
\(\sum_{i\in S}a_i\le s-1\). Modulo two the signs disappear, \(x(x-1)\) becomes
\(y=x^2+x\), \(x_j-1\) becomes \(1+x_j\), and by Lemma~\ref{lem:catalan} a term
survives exactly when every \(a_i\) is a power of two. So \(\widehat g_s\) reduces
to \(g_s\).

\begin{theorem}\label{thm:truncated-product}
For every finite \(\sigma\) with \(|\sigma|=n\) and every \(s\ge1\):
\begin{enumerate}
 \item \(\mult_a(g_s)\ge s\) for every nonzero \(a\in\Ft^\sigma\);
 \item \(g_s(\zero)=1\);
 \item \(\deg g_s\le n+2(s-1)-s_2(s-1)\).
\end{enumerate}
\lean{claims/TruncatedProduct.lean}
\end{theorem}

\begin{proof}
\emph{Expansion.} By Lemma~\ref{lem:telescope} with \(K=s\),
\[
 \prod_{i\in\sigma}(1+x_i)^{2^s}=\prod_{i\in\sigma}\Bigl((1+x_i)+\sum_{j<s}y_i^{2^j}\Bigr).
\]
Expanding the right-hand side gives the terms of Definition~\ref{def:truncated}
of every cost. Hence the left-hand side equals \(g_s+R\), where \(R\) collects the
terms of cost at least \(s\). This includes every term that uses a power
\(2^j>s-1\).

\emph{Multiplicity (1).}
\begin{itemize}
 \item Since \(u^2+u=0\) on \(\Ft\), each \(y_i\) vanishes at every point of
 \(\Ft^\sigma\). Lemma~\ref{lem:mult-basic}(3) then gives each term of \(R\)
 multiplicity at least its cost, so at least \(s\) everywhere. By
 Lemma~\ref{lem:mult-basic}(2), so does \(R\).
 \item At a nonzero \(a\), choose \(i_0\) with \(a_{i_0}=1\). Then
 \((1+x_{i_0})^{2^s}\) has multiplicity at least \(2^s\ge s\) at \(a\), and so does
 the whole product.
\end{itemize}
In characteristic two, \(g_s\) equals the product plus \(R\), so
\(\mult_a(g_s)\ge s\).

\emph{Value (2).} Every term with \(S\ne\emptyset\) vanishes at \(\zero\), and the term
with \(S=\emptyset\) is \(\prod(1+0)=1\).

\emph{Degree (3).} The term of \((S,p)\) has degree at most
\((n-|S|)+2\mu\), where \(\mu=\sum_{i\in S}p_i\le s-1\). By
Lemma~\ref{lem:digits}(1), \(|S|\ge s_2(\mu)\). By Lemma~\ref{lem:digits}(3),
\(2\mu-s_2(\mu)\le2(s-1)-s_2(s-1)\).
\end{proof}

\subsection{The hyperplane product}

Let \(F=\prod_{u\in\Ft^n\setminus\{\zero\}}(1+u\cdot x)\), where
\(u\cdot x=\sum_iu_ix_i\). Its \(2^n-1\) factors define the hyperplanes
\(u\cdot x=1\), which avoid the origin.

\begin{lemma}\label{lem:hyperplane-product}
For \(n\ge1\): \(\mult_a(F)\ge2^{n-1}\) for every nonzero \(a\in\Ft^n\),
\(F(\zero)=1\), and \(\deg F\le2^n-1\)
\lean{claims/PerOrderUpperBound.lean}.
\end{lemma}

\begin{proof}
Fix \(a\ne\zero\) and \(i\) with \(a_i=1\). Flipping \(u_i\) is an involution of
\(\Ft^n\) that changes \(u\cdot a\) by one, so exactly \(2^{n-1}\) vectors \(u\)
satisfy \(u\cdot a=1\). All of them are nonzero, so they index factors of \(F\)
that vanish at \(a\). Lemma~\ref{lem:mult-basic}(3,4) gives the first claim.
Each factor takes the value \(1\) at \(\zero\) and has degree at most one.
\end{proof}

These are the facts behind the hyperplane covers of~\cite{BBDM23}
(Remark~\ref{rem:hyperplane}).

\subsection{The construction}

\begin{theorem}\label{thm:construction}
Let \(\sigma\) be finite with \(|\sigma|=n\ge1\), fix \(x_1\in\sigma\), let
\(0\le\ell<k\), and write \(k-\ell-1=q2^n+r\) with \(0\le r<2^n\). The polynomial
\[
 P=y_1^{\ell}\,F^{2q}\,g_{r+1}
\]
satisfies:
\begin{enumerate}
 \item \(\mult_a(P)\ge k\) for every nonzero \(a\);
 \item \(\mult_{\zero}(P)=\ell\);
 \item \(\deg P\le\Phi(n,k,\ell)\).
\end{enumerate}
Hence \(\delta(n,k,\ell)\le\Phi(n,k,\ell)\)
\lean{claims/PerOrderUpperBound.lean}.
\end{theorem}

\begin{proof}
\begin{enumerate}
 \item Since \(u^2+u=0\) on \(\Ft\), \(y_1\) vanishes at every point, so
 \(\mult_a(y_1^\ell)\ge\ell\). With Lemma~\ref{lem:hyperplane-product},
 Theorem~\ref{thm:truncated-product}(1) and Lemma~\ref{lem:mult-basic}(3), at
 \(a\ne\zero\) we get
 \(\mult_a(P)\ge\ell+2q\cdot2^{n-1}+(r+1)=\ell+(k-\ell-1)+1=k\).
 \item \(y_1\) has origin order exactly one: its coefficient of \(x_1\) is \(1\).
 Also \(F(\zero)=g_{r+1}(\zero)=1\), so these factors have origin order zero.
 Multiplicities add over \(\Ft\) (Lemma~\ref{lem:mult-basic}(3)).
 \item \(\deg y_1^\ell\le2\ell\), \(\deg F^{2q}\le2q(2^n-1)\), and
 \(\deg g_{r+1}\le n+2r-s_2(r)\) by Theorem~\ref{thm:truncated-product}(3). The
 sum is \(n+2\ell+2(q2^n+r)-2q-s_2(r)=n+2k-2-2q-s_2(r)\), which is
 \(\Phi(n,k,\ell)\) by Lemma~\ref{lem:legendre}. \qedhere
\end{enumerate}
\end{proof}

For \(k-\ell-1<2^n\) we have \(q=0\), and \(P=y_1^\ell g_{k-\ell}\) is Menezes'
polynomial over \(\Ft\). Proposition~\ref{prop:menezes-polynomial} shows that it
works in every dimension; when \(k-\ell-1\ge2^n\) its degree bound exceeds
\(\Phi(n,k,\ell)\).

\begin{proposition}\label{prop:menezes-polynomial}
Let \(|\sigma|=n\ge1\), \(x_1\in\sigma\) and \(0\le\ell<k\). Then
\(P=y_1^\ell g_{k-\ell}\) has \(\mult_a(P)\ge k\) for every nonzero \(a\),
\(\mult_{\zero}(P)=\ell\), and \(\deg P\le n+2k-2-s_2(k-\ell-1)\)
\lean{claims/PerOrderConstruction.lean}.
\end{proposition}

\begin{proof}
As in the proof of Theorem~\ref{thm:construction}, with \(g_{k-\ell}\) in place of
\(F^{2q}g_{r+1}\): the multiplicity off the origin is at least \(\ell+(k-\ell)\),
the origin order is \(\ell\), and the degree is at most
\(2\ell+n+2(k-\ell-1)-s_2(k-\ell-1)\) by Theorem~\ref{thm:truncated-product}(3).
\end{proof}

Menezes proves the degree \(n+2k-2-s_2(k-\ell-1)\) optimal when \(n\ge k-1\)
\cite[Corollary~1.3]{Men26}. Section~\ref{sec:lower} proves
Theorem~\ref{thm:construction} optimal in every dimension; the construction has
degree exactly \(\Phi(n,k,\ell)\).
